\chapter{Capítulo 1}

\section{Características y objetivos}
\subsection{Compartición de recuros}
\begin{itemize}
    \item Recurso: Estos pueden ser tanto hardware como software. 
    \item Elasticidad: Además cabe la posibilidad de una gestión elástica 
    ,es decir, se pueden asignar mas recursos de forma momentánea.
    \item Políticas y métodos de gestión de recursos: Denominación ,Asociación y coordinación.
    \item Control descentralizado: Soporte a grupos de usarios que trabajan cooperativamente,
    a esto se les llama groupware
\end{itemize}
\subsection{Sistema abierto}
\begin{itemize}
    \item Determina si el sistema puede ser personaliza y ampliado
    \item Tanto software como hardware
    \item Unix es un ejemplo de sistema abierto
    \begin{itemize}
        \item Desarrolladores de aplicaciones pueden usar el sistema operativo sin necesidad de pagar por ello
        \item Vendedores de hardware, ya que el sistema puede ser ampliado
        \item Vendedores de software y usarios,ya que es independiente del hardware
    \end{itemize}
    \item Utilización de protocolos estándares de comunicación, tales como IPC.
\end{itemize}
\subsection{Concurrencia}
\begin{itemize}
    \item Concurrencia y paralelismo surgen en los sistemas distribuidos
    \begin{itemize}
        \item Actividad independiente de cada usuario
        \item Independencia entre recursos hardware
        \item Localización de los procesos en diferentes computadoras
    \end{itemize}
    \end{itemize}

\subsection{Escalabilidad}
\begin{itemize}
    \item Software debe ser independiente del tamaño.
    \item El proceso de petición de recursos debe ser independiente del tamaño de la red
    \item Desafiío, diseñar un software que sea escalable, eficiente y efectivo.Técnicas para conseguirlo:
    \begin{itemize}
        \item Utilización de cachés
        \item Replicación de datos
        \item Despliege de servicios
        \item Sistemas abiertos
    \end{itemize}
\end{itemize}

\subsection{Tolerancia a fallos}
\begin{itemize}
    \item El diseño se basa en dos propuestas:
    \begin{itemize}
        \item Redundancia hardware: Si un componente falla, otro puede tomar su lugar.
        \item Redundancia software: Si un proceso falla, otro puede tomar su lugar.
\end{itemize}
\end{itemize}
\subsection{Transparencia}
\begin{itemize}
    \item Ver el sistema como un todo
    \item Principial objetivo de diseño del software de sistemas distribuidos
    \item Tipos de transparencia básica:
    \begin{itemize}
        \item Transparencia de acceso: El usuario no debe preocuparse por la ubicación de los recursos.
        \item Transparencia de localización: El usuario no debe preocuparse por la ubicación física de los recursos.    
        \item Transperencia de concurrencia: Se puede dar operaciones concurrentes sobre los mismos recursos sin que interfieran entre sí.
    \end{itemize}
    \item Tipos de transparencia avanzada:
    \begin{itemize}
        \item Transparencia de replicación: El usuario no debe preocuparse por la existencia de múltiples copias de un recurso.
        \item Transparencia de fallo: Ocultación de fallos de hardware o software.
        \item Transparencia de migración: Migración de entidades sin afectas a programas o usuarios.
        \item Transparencia de rendimiento: El sistema puede adaptarse a las cargas de trabajo variables
        \item Transparencia de escalabilidad: El sistema puede crecer sin cambiar las estructuras de los sistemas.
    \end{itemize}
    
\end{itemize}

\section{Paradigmas de Aplicaciones Distribuidas}
\begin{itemize}


\item Aplicaciones paralelas de alto rendimiento:
    \begin{itemize}
        \item Decrementan el tiempo de respuesta
        \item Ventaja de escalabilidad frente a multiprocesadores
        \item Se clasifica en cuanto al grano de paralelismo:
        \begin{itemize}
            \item Grano fino: ++Tiempo Comunicacion, --Tiempo computación.
            \item Grano medio , adecuado para S.D y débilmente acoplados.
            \item Grano grueso: ++Tiempo computación, --Tiempo comunicacion.
    \end{itemize}
\item Aplicaciones tolerancia a fallos:
    \begin{itemize}
        \item Incrementan fiabilidad y disponibilidad
        \item Propiedad de fallo parcial
        \item Seguridad mediante la replicación en varios nodos
    \end{itemize}
\item Aplicaciones con especialización funcional:
    \begin{itemize}
        \item Cada aplicación es una colección de servicios
        \item Forma natural de diseñar estas aplicaciones es mediante un S.D con implementaciones alternativas (replicada,distribuidad,centralizada)
        \item Proporcionan alto rendimiento y seguridad
        \item Comunicación entre servicios
        \item Fácil escalabilidad de las aplicaciones
    \end{itemize}
\item Aplicaciones inherentemente distribuidas:
    \begin{itemize}
        \item Groupware: Sistema basado en computadores que soportan la colaboración entre usuarios.
        \begin{itemize}
            \item Ejemplos:
            \item Correo electrónico
            \item Software de gestión de flujos de trabajo
            \item Sistema Enterpirse Resource Planning (ERP)
            \item Business Process Management (BPMS)
        \end{itemize}
    \end{itemize}
\end{itemize}
\end{itemize}
\section{Características y objetivos emergentes en nuevos paradigmas de Computación Distribuida:}
\begin{itemize}
    \item Modularidad
    \item Cross platform(Interoperabilidad)
    \item Seguridad/privacidad
    \item Calidad de servicio
    \item Extensibilidad
    \item Movilidad
    \item Amigabilidad de las interfaces
    \item Elasticidad y escalabilidad
    \item Dispositovs móviles + LAN inalámbricas, sistemas distribuidos con clientes y tal vez servidores, móviles
    \item Permanecen los principios de diseño de sistemas distribuidos
    \item Aspectos adicionales a considerar:
    \begin{itemize}
        \item Variaciones impredecibles en la calidad de conexión de red
        \item Fiabilidad y robustez reducidad de los dispositivos móviles
        \item Potencias de los recursos reducida, derivada de las limitaciones de tamaño y peso
        \item Consumo de batería
    \end{itemize}
    \item Campo de investigación y desarrollo muy activo
    \item Áreas de trabajo:
    \begin{itemize}
        \item Redes móviles: IP, protocolos, TCP,UDP, etc.
        \item Acceso móvil a la información: Trabajo en modo desconexión,adaptación al ancho de banda, gestión de memoria.
        \item Soporte a aplicaciones adaptativas: Gestión de recursos adaptativa, transcoding.
        \item Técnicas de ahorro de energía: Gestión de energía adaptativa.
        \item Sensibilidad al contexto: Detección de la ubicación para determinar cambio en el comportamiento.
    \end{itemize}
    \item Invisibilidad
          \begin{itemize}
            \item Ideal expresado por Mark Weiser
            \item Minimizar las distracciones del usuario por el uso de la tecnología
            \item Si la tecnología no distrae, puede usarse a un nivel subconsciente
            \item Equilibrio frente a la ausencia total de notificaciones para evitar sorpresas posteriores
            \item Ejemplo: Conmutación transparente entre el uso de una red Wi-Fi y una conexión 4G.
          \end{itemize} 
    \item{Escalabilidad localizada}
    \begin{itemize}
        \item Al crecer el nivel de sofistificación de los espacios, aumentan las interacciones del usuario con el entorno.
        \item Consecuencias en el consumo de ancho de banda, energía y aparición de distractores.
        \item Escalabilidad
        \begin{itemize}
            \item Hasta ahora no importaba la distancia o el lugar del usuario.
            \item Ahora el número de interacciones debe reducirse en función de como se desplazan los usuarios.
        \end{itemize}
    \end{itemize}
    \item{Enmascaramiento de la condición desigual}
\end{itemize}
    



